\documentclass[10pt,a4paper]{amsart}
\usepackage{amsmath,amssymb,amsthm,graphicx}
% Retained minimal preamble for the Stacks Project chapter more-algebra.tex (More on Algebra).
% The chapter file pinned-original-source.tex carries only \input{preamble} on its first line, so no
% preamble text lives inside the chapter itself.  The Stacks Project uses one single shared preamble for
% every chapter; the snapshot stages it identically as sources/stacks-chow/preamble.tex and
% sources/stacks-etale-cohomology/preamble.tex.  The block below is copied mechanically, line for line,
% from sources/stacks-chow/preamble.tex lines 14-18 (the Xy-pic setup) and 154-217 (the native theorem
% declarations together with the mathematical macros).  Nothing was added, removed or redefined; the
% wrapper class and the amsmath/amssymb/amsthm/graphicx base packages are supplied by the assembler.
% Chapter: more-algebra.tex -- More on Algebra.
% For commutative diagrams we use Xy-pic
\usepackage[all]{xy}

% We use 2cell for 2-commutative diagrams.
\xyoption{2cell}

% Theorem environments.
%
\theoremstyle{plain}
\newtheorem{theorem}[subsection]{Theorem}
\newtheorem{proposition}[subsection]{Proposition}
\newtheorem{lemma}[subsection]{Lemma}

\theoremstyle{definition}
\newtheorem{definition}[subsection]{Definition}
\newtheorem{example}[subsection]{Example}
\newtheorem{exercise}[subsection]{Exercise}
\newtheorem{situation}[subsection]{Situation}

\theoremstyle{remark}
\newtheorem{remark}[subsection]{Remark}
\newtheorem{remarks}[subsection]{Remarks}

\numberwithin{equation}{subsection}

% Macros
%
\def\lim{\mathop{\mathrm{lim}}\nolimits}
\def\colim{\mathop{\mathrm{colim}}\nolimits}
\def\Spec{\mathop{\mathrm{Spec}}}
\def\Hom{\mathop{\mathrm{Hom}}\nolimits}
\def\Ext{\mathop{\mathrm{Ext}}\nolimits}
\def\SheafHom{\mathop{\mathcal{H}\!\mathit{om}}\nolimits}
\def\SheafExt{\mathop{\mathcal{E}\!\mathit{xt}}\nolimits}
\def\Sch{\mathit{Sch}}
\def\Mor{\mathop{\mathrm{Mor}}\nolimits}
\def\Ob{\mathop{\mathrm{Ob}}\nolimits}
\def\Sh{\mathop{\mathit{Sh}}\nolimits}
\def\NL{\mathop{N\!L}\nolimits}
\def\CH{\mathop{\mathrm{CH}}\nolimits}
\def\proetale{{pro\text{-}\acute{e}tale}}
\def\etale{{\acute{e}tale}}
\def\QCoh{\mathit{QCoh}}
\def\Ker{\mathop{\mathrm{Ker}}}
\def\Im{\mathop{\mathrm{Im}}}
\def\Coker{\mathop{\mathrm{Coker}}}
\def\Coim{\mathop{\mathrm{Coim}}}

% Boxtimes
%
\DeclareMathSymbol{\boxtimes}{\mathbin}{AMSa}{"02}

%
% Macros for moduli stacks/spaces
%
\def\QCohstack{\mathcal{QC}\!\mathit{oh}}
\def\Cohstack{\mathcal{C}\!\mathit{oh}}
\def\Spacesstack{\mathcal{S}\!\mathit{paces}}
\def\Quotfunctor{\mathrm{Quot}}
\def\Hilbfunctor{\mathrm{Hilb}}
\def\Curvesstack{\mathcal{C}\!\mathit{urves}}
\def\Polarizedstack{\mathcal{P}\!\mathit{olarized}}
\def\Complexesstack{\mathcal{C}\!\mathit{omplexes}}
% \Pic is the operator that assigns to X its picard group, usage \Pic(X)
% \Picardstack_{X/B} denotes the Picard stack of X over B
% \Picardfunctor_{X/B} denotes the Picard functor of X over B
\def\Pic{\mathop{\mathrm{Pic}}\nolimits}
\def\Picardstack{\mathcal{P}\!\mathit{ic}}
\def\Picardfunctor{\mathrm{Pic}}
\def\Deformationcategory{\mathcal{D}\!\mathit{ef}}
\usepackage{cleveref}
\crefname{originallemma}{original lemma}{original lemmas}
\Crefname{originallemma}{Original lemma}{Original lemmas}
\makeatletter
\newlabel{more-algebra-lemma-torsion-free-finite-noetherian-domain}{{0AUV}{source}{}{sourcepointer-1}{}}
\newlabel{more-algebra-lemma-torsion-free-finite-noetherian-domain@cref}{{[originallemma][0][]0AUV}{[1][0][]source}{}{}{}}
\makeatother
\crefname{originallemma}{original lemma}{original lemmas}
\Crefname{originallemma}{Original lemma}{Original lemmas}
\makeatletter
\newlabel{lemma-finite-torsion-free-submodule-free}{{0AUU}{source}{}{sourcepointer-2}{}}
\newlabel{lemma-finite-torsion-free-submodule-free@cref}{{[originallemma][0][]0AUU}{[1][0][]source}{}{}{}}
\makeatother
\crefname{originallemma}{original lemma}{original lemmas}
\Crefname{originallemma}{Original lemma}{Original lemmas}
\makeatletter
\newlabel{lemma-localize-torsion}{{0AUR}{source}{}{sourcepointer-3}{}}
\newlabel{lemma-localize-torsion@cref}{{[originallemma][0][]0AUR}{[1][0][]source}{}{}{}}
\makeatother
\crefname{originallemma}{original lemma}{original lemmas}
\Crefname{originallemma}{Original lemma}{Original lemmas}
\makeatletter
\newlabel{algebra-lemma-criterion-no-embedded-primes}{{031Q}{source}{}{sourcepointer-4}{}}
\newlabel{algebra-lemma-criterion-no-embedded-primes@cref}{{[originallemma][0][]031Q}{[1][0][]source}{}{}{}}
\makeatother
\crefname{originallemma}{original lemma}{original lemmas}
\Crefname{originallemma}{Original lemma}{Original lemmas}
\makeatletter
\newlabel{algebra-lemma-ass}{{02M3}{source}{}{sourcepointer-5}{}}
\newlabel{algebra-lemma-ass@cref}{{[originallemma][0][]02M3}{[1][0][]source}{}{}{}}
\makeatother
\crefname{originallemma}{original lemma}{original lemmas}
\Crefname{originallemma}{Original lemma}{Original lemmas}
\makeatletter
\newlabel{algebra-lemma-ass-zero}{{0587}{source}{}{sourcepointer-6}{}}
\newlabel{algebra-lemma-ass-zero@cref}{{[originallemma][0][]0587}{[1][0][]source}{}{}{}}
\makeatother
\title{Five equivalent characterisations of torsion freeness for a nonzero finite module over a Noetherian domain}
\author{The Stacks Project Authors}
\date{stacks-project:more-algebra.tex:local-snapshot-20261012; GNU FDL 1.2 or later}
\begin{document}
\maketitle
\setcounter{subsection}{1}
\noindent\textit{Editorial note.} All delivered mathematical text is the author's own, copied from the pinned Stacks Project chapter file more-algebra.tex; no statement or proof step was rewritten, reordered or removed.\par\medskip
\noindent\textit{Editorial note.} The chapter file pinned-original-source.tex carries only \textbackslash{}input\{preamble\} on its first line, so no preamble text lives inside the chapter. The retained file sources/stacks-more-algebra/retained-preamble.tex is a mechanical, line-for-line copy of the project's single shared preamble as staged in the snapshot, namely sources/stacks-chow/preamble.tex lines 14-18 (Xy-pic setup) and 154-217 (theorem declarations and mathematical macros). Nothing was added, removed or redefined, no package was added, and the excerpt serves only as preamble, never as a span coordinate.\par\medskip
\noindent\textit{Editorial note.} Both counted spans start at the \textbackslash{}begin line of their environment, so the statement and the author's proof are delivered whole; surrounding prose of the section is deliberately outside the counted spans and is carried as standard prerequisites instead.\par\medskip
\noindent\textit{Editorial note.} The citation fields (source\_id, source\_version, authors, source\_url, license.name, source\_title) are identical to those of the cards already published from this chapter, so that the batch selector treats them as one source instead of deferring them as a same-source clash.\par\medskip
\noindent\textit{Allowed standard prerequisites (editorial).}\begin{itemize}
\item Torsion element and torsion free module over a domain, and the equivalence with being a submodule of a finite free module. definition at the start of the section, source lines 4781-4792, and the statement counted by the sibling card S1-algebra-finite-torsion-free-submodule-free at 4863-4881. Supplies the first two conditions and the equivalence between them, which is quoted at the start of the proof.
\item Associated primes of a finite module over a Noetherian ring, property (S\_1), and the criterion for the absence of embedded associated primes. the author cites the Algebra chapter lemma algebra-lemma-criterion-no-embedded-primes (tag 031Q) at proof lines 4899-4901. Gives the equivalence of conditions (4) and (5), the only input coming from outside this chapter.
\item A torsion element produces an embedding of R/I for a nonzero ideal I, hence an associated prime different from the zero ideal. the author cites Algebra lemmas algebra-lemma-ass (tag 02M3) and algebra-lemma-ass-zero (tag 0587) at proof lines 4906-4909. Provides the contradiction that finishes the implication (5) implies (1).
\end{itemize}\medskip
\ifdefined\hypertarget\hypertarget{sourcepointer-1}{}\fi
\noindent\textit{Source reference.} Original lemma 0AUV: focused-seeds/AGP/sources/stacks-snapshot-20261011/tags.txt, source lines 13663--13663. The lemma's own canonical tag, quoted so that readers can resolve the label inside the delivered text; it points at the counted statement itself and carries no proof content.\par\medskip
\ifdefined\hypertarget\hypertarget{sourcepointer-2}{}\fi
\noindent\textit{Source reference.} Original lemma 0AUU: focused-seeds/AGP/sources/stacks-snapshot-20261011/tags.txt, source lines 13662--13662. Navigation pointer to the preceding statement of the same section, which supplies the equivalence of conditions (1) and (2); that statement is counted by the sibling card S1-algebra-finite-torsion-free-submodule-free of this lane and nothing here re-proves it.\par\medskip
\ifdefined\hypertarget\hypertarget{sourcepointer-3}{}\fi
\noindent\textit{Source reference.} Original lemma 0AUR: focused-seeds/AGP/sources/stacks-snapshot-20261011/tags.txt, source lines 13659--13659. Navigation pointer to the statement of the same section that a localisation of a torsion free module is torsion free; it is used at one step of the proof, whose remaining implications are carried out inside the counted span.\par\medskip
\ifdefined\hypertarget\hypertarget{sourcepointer-4}{}\fi
\noindent\textit{Source reference.} Original lemma 031Q: focused-seeds/AGP/sources/stacks-snapshot-20261011/tags.txt, source lines 3993--3993. Navigation pointer into the Algebra chapter for the criterion identifying the conditions (4) and (5) on associated primes; it is a standard input, not part of the argument carried here.\par\medskip
\ifdefined\hypertarget\hypertarget{sourcepointer-5}{}\fi
\noindent\textit{Source reference.} Original lemma 02M3: focused-seeds/AGP/sources/stacks-snapshot-20261011/tags.txt, source lines 3430--3430. Navigation pointer into the Algebra chapter for the connection between embeddings of quotients by an ideal and associated primes, used in the last step of the proof.\par\medskip
\ifdefined\hypertarget\hypertarget{sourcepointer-6}{}\fi
\noindent\textit{Source reference.} Original lemma 0587: focused-seeds/AGP/sources/stacks-snapshot-20261011/tags.txt, source lines 6770--6770. Navigation pointer into the Algebra chapter for the behaviour of associated primes with respect to the zero ideal, used together with the previous pointer to conclude that the prime found is embedded.\par\medskip
\begin{lemma}
\label{lemma-torsion-free-finite-noetherian-domain}
Let $R$ be a Noetherian domain. Let $M$ be a nonzero finite $R$-module.
The following are equivalent
\begin{enumerate}
\item $M$ is torsion free,
\item $M$ is a submodule of a finite free module,
\item $(0)$ is the only associated prime of $M$,
\item $(0)$ is in the support of $M$ and $M$ has property $(S_1)$, and
\item $(0)$ is in the support of $M$ and $M$ has no embedded associated prime.
\end{enumerate}
\end{lemma}

\begin{proof}
We have seen the equivalence of (1) and (2) in
Lemma \ref{lemma-finite-torsion-free-submodule-free}.
We have seen the equivalence of (4) and (5) in
Algebra, Lemma \ref{algebra-lemma-criterion-no-embedded-primes}.
The equivalence between (3) and (5) is immediate from the definition.
A localization of a torsion free module is torsion free
(Lemma \ref{lemma-localize-torsion}), hence it is clear that a
$M$ has no associated primes different from $(0)$. Thus (1)
implies (5). Conversely, assume (5). If $M$ has torsion,
then there exists an embedding $R/I \subset M$ for some nonzero
ideal $I$ of $R$. Hence $M$ has an associated prime different
from $(0)$
(see Algebra, Lemmas \ref{algebra-lemma-ass} and \ref{algebra-lemma-ass-zero}).
This is an embedded associated prime which contradicts the assumption.
\end{proof}

\end{document}
